% ============================================================== 6. Conclusion
\section{Conclusion}
\label{sec:conclusion}

The Fed-ISIC2019 benchmark split places images of the same lesion in training and test sets, unevenly
across its acquisition-source clients. Models look about 8 balanced-accuracy points better on the leaky
test sets than on their new-lesion subsets, and up to 16 points better for centralized training.
In a randomized experiment on fixed test images, which covers lesions with several images, adding sibling
images to training raised balanced accuracy by 9.1 points for FedAvg, with effects between 4.5 and 28.8
points across training regimes. Leakage also changed which methods appeared best, mainly among classifiers
that differed in head architecture or training regime. At the level of individual clients, a shared conformal threshold under-covered some
sites; an empirical-Bayes shrinkage of per-client thresholds kept mean client coverage near the target and
improved the worst-client tail for one of two encoders, while a fixed shrinkage rule failed its
pre-specified test. On unseen clients and smartphone images, reliability dropped further, and local
recalibration restored coverage only with large prediction sets. We recommend that federated medical
imaging benchmarks with repeated imaging use lesion- or patient-grouped splits and report calibration and
coverage per client, and we release the corrected splits, code, predictions and protocol history to
support this.
